\documentclass{article}
% notation.tex --- shared preamble for Warp paper-mods.
%
% Every mod*.tex file begins with:
%
%     % notation.tex --- shared preamble for Warp paper-mods.
%
% Every mod*.tex file begins with:
%
%     % notation.tex --- shared preamble for Warp paper-mods.
%
% Every mod*.tex file begins with:
%
%     \input{notation.tex}
%
% Conventions
% -----------
% * Framing: IOP in the sense of Ben-Sasson--Chiesa--Spooner (BCS, TCC 2016),
%   NOT Algebraic Holographic Proofs. Oracles are functions $f:[n] \to \F$,
%   queried by index; multilinear-extension semantics layer on top for point
%   queries.
% * Every modification file cites the Rust module path it pairs with.

\usepackage{amsmath,amssymb,amsthm}
\usepackage{mathtools}
\usepackage{hyperref}

% ---- Fields, indexing, domains -------------------------------------------
\newcommand{\F}{\mathbb{F}}
\newcommand{\idx}[1]{\ensuremath{[#1]}}               % index set [n]
\newcommand{\bool}[1]{\ensuremath{\{0,1\}^{#1}}}      % boolean hypercube
\newcommand{\size}[1]{\ensuremath{\left\lvert#1\right\rvert}}

% ---- Oracles (Modification 1) --------------------------------------------
% An Oracle carries three views of the same object:
%   (i)   an evaluation table  f : [n] -> F        (BCS-native)
%   (ii)  a multilinear extension  \hat f : F^{log n} -> F
%   (iii) a commitment handle (Merkle root under BCS compilation)
% See mod1_oracle.tex.
\newcommand{\Oracle}[1]{\ensuremath{\mathsf{O}[#1]}}
\newcommand{\MLE}[1]{\ensuremath{\widehat{#1}}}
\newcommand{\Commit}[1]{\ensuremath{\mathsf{com}(#1)}}
\newcommand{\query}{\ensuremath{\mathsf{query}}}      % generic query operator

% ---- Accumulator state (Modification 3, reserved) -------------------------
\newcommand{\AccX}{\ensuremath{\mathsf{acc}_x}}       % accumulator instance
\newcommand{\AccW}{\ensuremath{\mathsf{acc}_w}}       % accumulator witness
\newcommand{\AccState}{\ensuremath{(\AccX, \AccW)}}

% ---- Common structural IORs ----------------------------------------------
\newcommand{\IOR}[1]{\textsc{#1}}
\newcommand{\Zerocheck}{\IOR{Zerocheck}}
\newcommand{\Sumcheck}{\IOR{Sumcheck}}
\newcommand{\Pesat}{\IOR{Pesat}}
\newcommand{\TwinConstraint}{\IOR{TwinConstraint}}
\newcommand{\Proximity}{\IOR{Proximity}}
\newcommand{\Batching}{\IOR{Batching}}
\newcommand{\OOD}{\IOR{OOD}}

% ---- Transcript notation --------------------------------------------------
\newcommand{\absorb}{\ensuremath{\lhd}}               % absorb into transcript
\newcommand{\squeeze}{\ensuremath{\rhd}}              % derive from transcript

% ---- Code / module cross-reference ---------------------------------------
% Use \codemod{src/protocol/oracle.rs} in .tex files to point at a Rust
% module. Renders as inline typewriter with a URL when compiled with hyperref.
\newcommand{\codemod}[1]{\href{../../#1}{\texttt{#1}}}

%
% Conventions
% -----------
% * Framing: IOP in the sense of Ben-Sasson--Chiesa--Spooner (BCS, TCC 2016),
%   NOT Algebraic Holographic Proofs. Oracles are functions $f:[n] \to \F$,
%   queried by index; multilinear-extension semantics layer on top for point
%   queries.
% * Every modification file cites the Rust module path it pairs with.

\usepackage{amsmath,amssymb,amsthm}
\usepackage{mathtools}
\usepackage{hyperref}

% ---- Fields, indexing, domains -------------------------------------------
\newcommand{\F}{\mathbb{F}}
\newcommand{\idx}[1]{\ensuremath{[#1]}}               % index set [n]
\newcommand{\bool}[1]{\ensuremath{\{0,1\}^{#1}}}      % boolean hypercube
\newcommand{\size}[1]{\ensuremath{\left\lvert#1\right\rvert}}

% ---- Oracles (Modification 1) --------------------------------------------
% An Oracle carries three views of the same object:
%   (i)   an evaluation table  f : [n] -> F        (BCS-native)
%   (ii)  a multilinear extension  \hat f : F^{log n} -> F
%   (iii) a commitment handle (Merkle root under BCS compilation)
% See mod1_oracle.tex.
\newcommand{\Oracle}[1]{\ensuremath{\mathsf{O}[#1]}}
\newcommand{\MLE}[1]{\ensuremath{\widehat{#1}}}
\newcommand{\Commit}[1]{\ensuremath{\mathsf{com}(#1)}}
\newcommand{\query}{\ensuremath{\mathsf{query}}}      % generic query operator

% ---- Accumulator state (Modification 3, reserved) -------------------------
\newcommand{\AccX}{\ensuremath{\mathsf{acc}_x}}       % accumulator instance
\newcommand{\AccW}{\ensuremath{\mathsf{acc}_w}}       % accumulator witness
\newcommand{\AccState}{\ensuremath{(\AccX, \AccW)}}

% ---- Common structural IORs ----------------------------------------------
\newcommand{\IOR}[1]{\textsc{#1}}
\newcommand{\Zerocheck}{\IOR{Zerocheck}}
\newcommand{\Sumcheck}{\IOR{Sumcheck}}
\newcommand{\Pesat}{\IOR{Pesat}}
\newcommand{\TwinConstraint}{\IOR{TwinConstraint}}
\newcommand{\Proximity}{\IOR{Proximity}}
\newcommand{\Batching}{\IOR{Batching}}
\newcommand{\OOD}{\IOR{OOD}}

% ---- Transcript notation --------------------------------------------------
\newcommand{\absorb}{\ensuremath{\lhd}}               % absorb into transcript
\newcommand{\squeeze}{\ensuremath{\rhd}}              % derive from transcript

% ---- Code / module cross-reference ---------------------------------------
% Use \codemod{src/protocol/oracle.rs} in .tex files to point at a Rust
% module. Renders as inline typewriter with a URL when compiled with hyperref.
\newcommand{\codemod}[1]{\href{../../#1}{\texttt{#1}}}

%
% Conventions
% -----------
% * Framing: IOP in the sense of Ben-Sasson--Chiesa--Spooner (BCS, TCC 2016),
%   NOT Algebraic Holographic Proofs. Oracles are functions $f:[n] \to \F$,
%   queried by index; multilinear-extension semantics layer on top for point
%   queries.
% * Every modification file cites the Rust module path it pairs with.

\usepackage{amsmath,amssymb,amsthm}
\usepackage{mathtools}
\usepackage{hyperref}

% ---- Fields, indexing, domains -------------------------------------------
\newcommand{\F}{\mathbb{F}}
\newcommand{\idx}[1]{\ensuremath{[#1]}}               % index set [n]
\newcommand{\bool}[1]{\ensuremath{\{0,1\}^{#1}}}      % boolean hypercube
\newcommand{\size}[1]{\ensuremath{\left\lvert#1\right\rvert}}

% ---- Oracles (Modification 1) --------------------------------------------
% An Oracle carries three views of the same object:
%   (i)   an evaluation table  f : [n] -> F        (BCS-native)
%   (ii)  a multilinear extension  \hat f : F^{log n} -> F
%   (iii) a commitment handle (Merkle root under BCS compilation)
% See mod1_oracle.tex.
\newcommand{\Oracle}[1]{\ensuremath{\mathsf{O}[#1]}}
\newcommand{\MLE}[1]{\ensuremath{\widehat{#1}}}
\newcommand{\Commit}[1]{\ensuremath{\mathsf{com}(#1)}}
\newcommand{\query}{\ensuremath{\mathsf{query}}}      % generic query operator

% ---- Accumulator state (Modification 3, reserved) -------------------------
\newcommand{\AccX}{\ensuremath{\mathsf{acc}_x}}       % accumulator instance
\newcommand{\AccW}{\ensuremath{\mathsf{acc}_w}}       % accumulator witness
\newcommand{\AccState}{\ensuremath{(\AccX, \AccW)}}

% ---- Common structural IORs ----------------------------------------------
\newcommand{\IOR}[1]{\textsc{#1}}
\newcommand{\Zerocheck}{\IOR{Zerocheck}}
\newcommand{\Sumcheck}{\IOR{Sumcheck}}
\newcommand{\Pesat}{\IOR{Pesat}}
\newcommand{\TwinConstraint}{\IOR{TwinConstraint}}
\newcommand{\Proximity}{\IOR{Proximity}}
\newcommand{\Batching}{\IOR{Batching}}
\newcommand{\OOD}{\IOR{OOD}}

% ---- Transcript notation --------------------------------------------------
\newcommand{\absorb}{\ensuremath{\lhd}}               % absorb into transcript
\newcommand{\squeeze}{\ensuremath{\rhd}}              % derive from transcript

% ---- Code / module cross-reference ---------------------------------------
% Use \codemod{src/protocol/oracle.rs} in .tex files to point at a Rust
% module. Renders as inline typewriter with a URL when compiled with hyperref.
\newcommand{\codemod}[1]{\href{../../#1}{\texttt{#1}}}


\title{Modification 1: \\ Oracle as first-class IOR output}
\author{Warp paper-mods}
\date{}

\begin{document}
\maketitle

\paragraph{Paired Rust module.} \codemod{src/protocol/oracle.rs}, consumed
throughout \codemod{src/protocol/phases/}.

\section{Motivation}

The Warp paper decomposes its construction as a sequence of interactive
oracle reductions (IORs). In the pure IOR formalism, each IOR's prover is an
abstract Turing machine that recomputes anything it needs; there is no notion
of ``shared state'' across IORs beyond the transcript. This abstraction is
sound but imposes an unpleasant cost on any implementation: several distinct
IORs in Warp --- the twin-constraint sumcheck, OOD sampling, the batching
sumcheck, and the proximity-query phase --- all operate on \emph{the same
codeword} \(f\), its multilinear extension \(\MLE{f}\), and its Merkle
commitment. A literal reading of the paper forces the code either to
rematerialise these objects phase by phase (wasteful) or to smuggle them
across phase boundaries as ambient state (unprincipled).

The BCS formalism already provides the right object to legitimise the
sharing: an \emph{oracle message}. BCS treats prover oracles as typed,
persistent objects that can be queried in any subsequent round. What the
paper presently leaves implicit is that \emph{composed} IORs can export and
import these oracles as typed inputs and outputs.

This modification promotes Oracle to a first-class input/output type of every
Warp IOR. No new soundness argument is required; the BCS compiler already
handles the Merkle-commitment bookkeeping. What changes is the presentation,
which in turn legitimises a direct Rust implementation in
\codemod{src/protocol/oracle.rs}.

\section{Definition}

An \emph{oracle} \(\Oracle{f}\) carries the following views of a single
object:
\begin{itemize}
  \item \(f : \idx{n} \to \F\) is an evaluation table (the \emph{BCS-native}
        view);
  \item \(\MLE{f} : \F^{\log n} \to \F\) is the unique multilinear polynomial
        over \(\bool{\log n}\) agreeing with \(f\) after identifying
        \(\idx{n} \leftrightarrow \bool{\log n}\) by binary expansion;
  \item under BCS compilation, \(f\) is committed via a Merkle root
        \(\Commit{f}\); this commitment is produced by a separate
        \emph{commit} operation and is not a field of \(\Oracle{f}\).
\end{itemize}

\(\MLE{f}\) is \emph{implied} by \(f\); the prover materialises it lazily and
caches. See \codemod{src/protocol/oracle.rs}, method \texttt{query\_at\_point}.

\paragraph{Implementation note: commitments and oracles are not 1:1.} Warp's
\Pesat{} phase interleaves \(l_1\) codewords into a single Merkle tree
(\codemod{src/crypto/merkle/mod.rs}, \texttt{build\_codeword\_leaves}), so one
commitment covers many oracles. The Rust \codemod{Oracle} struct therefore
stores only \(f\) and the lazy \(\MLE{f}\); the Merkle tree implementing
\(\Commit{f}\) is tracked by the enclosing data structure
(\codemod{src/types.rs} --- \texttt{PesatOutput}, \texttt{AccumulatorWitness}).
This is an orthogonal implementation choice; the IOR composition rule in
\S4 is unchanged.

\section{Query interfaces}

Two legal query operations on \(\Oracle{f}\):

\begin{description}
  \item[Index query.] \(f[i]\) for \(i \in \idx{n}\). BCS-native.
        The verifier receives \(f[i]\) alongside a Merkle opening against
        \(\Commit{f}\); the BCS compiler makes this non-interactive via
        Fiat--Shamir. Corresponds to
        \codemod{src/protocol/oracle.rs}, method \texttt{query\_at\_leaf}.
  \item[Point query.] \(\MLE{f}(\zeta)\) for \(\zeta \in \F^{\log n}\).
        The prover answers directly; the verifier's check that the answer is
        consistent with \(\Commit{f}\) is deferred to a downstream IOR (in
        Warp, the batching sumcheck). Corresponds to
        \codemod{src/protocol/oracle.rs}, method \texttt{query\_at\_point}.
\end{description}

Both operations act on \emph{the same} oracle; a downstream IOR is free to
mix query kinds.

\section{Composition rule}

If IOR \(A\) has signature
\(\; (\text{stmt}_A, \text{wit}_A) \;\mapsto\; (\text{stmt}_A', \Oracle{f})\)
and IOR \(B\) has signature
\(\; (\text{stmt}_B, \Oracle{f}) \;\mapsto\; \text{stmt}_B'\),
then the sequential composition \(A; B\) has signature
\(\; (\text{stmt}_A, \text{stmt}_B, \text{wit}_A) \;\mapsto\; (\text{stmt}_A', \text{stmt}_B') \).
Soundness follows from the BCS composition theorem for oracle messages: the
Merkle commitment \(\Commit{f}\) is produced once by \(A\) and reused by any
number of downstream IORs.

\section{Application to Warp}

Every Warp phase can be restated in terms of oracles:

\begin{center}
\begin{tabular}{l l}
\toprule
Phase & Oracle role \\
\midrule
Encode-and-commit   & emits \(\Oracle{f}\) per fresh witness \\
\Pesat              & emits \(\Oracle{u}\) for accumulated + fresh codewords \\
\TwinConstraint     & consumes \(\Oracle{u}\), emits reduced claim \\
\OOD                & point queries on \(\Oracle{f}\) \\
\Batching           & consumes \(\Oracle{f}\); emits new \(\Oracle{f}'\) via fold \\
\Proximity          & index queries on \(\Oracle{f}\) with auth paths \\
\bottomrule
\end{tabular}
\end{center}

The same object carries the codeword, the multilinear extension, and the
Merkle tree; downstream phases access whichever view the operation requires.
In Rust this is literally one struct (\codemod{src/protocol/oracle.rs}).

\section{What this does not change}

Soundness proofs, parameter choices, and the transcript layout are unchanged.
Modification 1 is presentational: it makes the existing data flow type-check
in a paper that stays inside BCS rather than retreating to AHP.

\end{document}
